% year: תשפ"ה
% type: בוחן
% semester: א
% source: exams/בחנים/תשפ״ה, בוחן, 9141.pdf
% part: ב
% number: 7
% points: 10
% categories: continuity, func-limits, multiple-choice

\begin{question}
בדקו רציפות של הפונקציה הבאה בנקודה 2:
\[
f(x)=\begin{cases}
\frac{1}{1+4^{\frac{1}{x-2}}}, & \text{if } x\neq 2,\\
0, & \text{if } x=2
\end{cases}
\]
אילו מהטענות הבאות נכונה?
\begin{enumerate}
  \item הפונקציה רציפה ב-2.
  \item הפונקציה בעלת אי-רציפות סליקה ב-2.
  \item הפונקציה בעלת אי-רציפות מסוג קפיצה ב-2.
  \item הפונקציה בעלת אי-רציפות עיקרית ב-2.
\end{enumerate}
\end{question}

\begin{hint}
חשבו בנפרד את הגבולות החד-צדדיים: לאן שואף $\frac{1}{x-2}$ כאשר $x\to2^+$ וכאשר $x\to2^-$?
\end{hint}

\begin{solution}
\textbf{התשובה הנכונה: (ג)}.

נחשב את הגבולות החד-צדדיים ב-2.

\textbf{מימין:} כאשר $x\to2^+$, $x-2\to0^+$ ולכן $\frac1{x-2}\to+\infty$. מכיוון ש-$4>1$, $\lim_{t\to+\infty}4^t=\infty$, ולכן $1+4^{\frac{1}{x-2}}\to\infty$ ו-
\[
\lim_{x\to2^+}f(x)=0 .
\]
\textbf{משמאל:} כאשר $x\to2^-$, $\frac1{x-2}\to-\infty$ ו-$\lim_{t\to-\infty}4^t=0$, ולכן
\[
\lim_{x\to2^-}f(x)=\frac{1}{1+0}=1 .
\]
שני הגבולות החד-צדדיים קיימים (סופיים) אך שונים, ולכן הגבול $\lim_{x\to2}f(x)$ אינו קיים, וזו בדיוק אי-רציפות מסוג קפיצה (מסוג ראשון לא סליקה).

\textbf{למה האחרות שגויות:} (א) — הגבול ב-2 לא קיים, ולכן $f$ אינה רציפה ב-2 (העובדה ש-$f(2)=0$ שווה לגבול מימין נותנת רק רציפות מימין). (ב) — אי-רציפות סליקה דורשת שהגבול הדו-צדדי יהיה קיים. (ד) — אי-רציפות עיקרית (מסוג שני) פירושה שלפחות אחד הגבולות החד-צדדיים אינו קיים כגבול סופי, וכאן שניהם קיימים.
\end{solution}
